\chapter{Consumo energetico}
\label{chap:consumi}

Questo capitolo stima il consumo energetico dei sensori e del nodo che li ospita. L'analisi è a titolo informativo e usa i valori dichiarati dai produttori.

Il capitolo è organizzato come segue. La Sezione~\ref{sec:consumi-dati} riprende i consumi dei sensori dal Capitolo~\ref{chap:moduli} e vi aggiunge quelli del microcontrollore. La Sezione~\ref{sec:consumi-nodi} li somma nelle configurazioni del nodo, e la Sezione~\ref{sec:autonomia} ne calcola l'autonomia a batteria. La Sezione~\ref{sec:consumi-vincoli} ne trae i vincoli per DIPME.

\section{Dati di partenza}
\label{sec:consumi-dati}

Le correnti dei sensori sono riportate nelle specifiche del Capitolo~\ref{chap:moduli}, e per i conti di questo capitolo basta la corrente media di ciascuno. L'\ldb\ assorbe in media 80~mA a 5~V, pari a circa 400~mW, e l'\ldc\ 50~mA a 3,3~V, circa 165~mW. L'\pir\ a riposo resta sotto i 50~\textmu A, cioè circa 0,3~mW. Per l'ESP32-WROOM-32 le correnti nelle diverse modalità operative sono riportate nella Tabella~\ref{tab:consumi-esp32}, dal datasheet Espressif~\cite{esp32datasheet}.

\begin{table}[htbp]
\centering
\caption{Consumo dell'ESP32-WROOM-32 per modalità operativa.}
\label{tab:consumi-esp32}
\small
\begin{tabular}{@{}lr@{}}
\toprule
\textbf{Modalità} & \textbf{Corrente} \\
\midrule
Attivo, WiFi in trasmissione (picco) & 160--260~mA \\
Attivo, WiFi connesso in ascolto & 95--100~mA \\
Modem-sleep (CPU attiva a 80--240~MHz, WiFi spento) & 20--68~mA \\
Light-sleep & 0,8~mA \\
Deep-sleep (solo dominio RTC) & 0,01~mA \\
\bottomrule
\end{tabular}
\end{table}

Il nodo del progetto DIPME usa invece un microcontrollore STM32WLE con radio LoRa e un PIR digitale Panasonic EKMB139~\cite{callisto2026dipme}, componenti a bassissimo consumo pensati per la batteria. I valori di questo capitolo si riferiscono al kit degli esperimenti, ma anche con componenti diversi è il radar a determinare il consumo finale del nodo.

\section{Il nodo completo}
\label{sec:consumi-nodi}

Sommando microcontrollore e sensore si ottiene il consumo del nodo completo (Tabella~\ref{tab:consumi-nodi}). Le prime tre righe sono le configurazioni usate negli esperimenti, le ultime due quelle con il solo PIR, con il microcontrollore acceso oppure in deep-sleep.

\begin{table}[htbp]
\centering
\caption{Consumo stimato del nodo completo nelle configurazioni di interesse, somma
dei valori di datasheet.}
\label{tab:consumi-nodi}
\small
\begin{tabular}{@{}lrr@{}}
\toprule
\textbf{Configurazione} & \textbf{Corrente} & \textbf{Potenza a 5~V} \\
\midrule
ESP32 + \ldb\ + WiFi (web UI) & 175--340~mA & 0,9--1,7~W \\
ESP32 + \ldb, WiFi spento (logger seriale) & 100--150~mA & 0,5--0,75~W \\
ESP32 + \ldc, WiFi spento & 70--120~mA & 0,35--0,6~W \\
ESP32 + PIR, WiFi spento & 20--70~mA & 0,1--0,35~W \\
ESP32 in deep-sleep + PIR, sveglia su interrupt & $\mathbf{0{,}06}$~\textbf{mA} & 0,3~mW \\
\bottomrule
\end{tabular}
\end{table}

L'ultima riga è il punto chiave, perché l'uscita digitale del PIR può risvegliare
l'ESP32 dal deep-sleep con un interrupt su GPIO. Il nodo può quindi restare in attesa a
costo praticamente nullo, con il PIR nel ruolo di guardiano.

\section{Autonomia a batteria}
\label{sec:autonomia}

Una cella litio-ione 18650 da 3000~mAh, con un regolatore all'85\,\% di rendimento, fornisce circa 2550~mAh utili, che divisi per la corrente media danno le autonomie della Tabella~\ref{tab:autonomia} e della Figura~\ref{fig:consumi}. Il calcolo non include l'autoscarica della cella.

\begin{table}[htbp]
\centering
\caption{Autonomia stimata con una cella 18650 da 3000~mAh.}
\label{tab:autonomia}
\small
\begin{tabular}{@{}lrr@{}}
\toprule
\textbf{Scenario} & \textbf{Corrente media} & \textbf{Autonomia} \\
\midrule
ESP32 + \ldb\ + WiFi (web UI) & 200~mA & 13~h \\
ESP32 + \ldb, WiFi spento & 130~mA & 20~h \\
ESP32 + PIR, WiFi spento & 25~mA & 4 giorni \\
ESP32 in deep-sleep + PIR & 0,06~mA & 4,9 anni \\
\bottomrule
\end{tabular}
\end{table}

\begin{figure}[htbp]
\centering
\includegraphics[width=0.9\linewidth]{Consumi}
\caption{Consumi dei componenti e autonomia dei nodi, in scala logaritmica.}
\label{fig:consumi}
\end{figure}

\section{Cosa indicano i consumi}
\label{sec:consumi-vincoli}

Con il radar sempre acceso la batteria dura meno di un giorno (Tabella~\ref{tab:autonomia}), mentre un nodo in deep-sleep svegliato dal PIR arriva a 4,9 anni, un limite superiore perché durante le lezioni il PIR lo sveglia di continuo. Accendendo il radar un minuto ogni cinque, con l'ESP32 in light-sleep negli altri quattro, la corrente media scende a circa 27~mA e l'autonomia supera i tre giorni, con una latenza fino a quattro minuti. Su questi valori si basa l'architettura delle conclusioni (Capitolo~\ref{chap:conclusioni}), in cui il PIR decide quando accendere il radar.
